% 2010 年天津卷理科第 12 题：正视图、侧视图与俯视图。
% 来源 PDF：普通高考/2010/2010天津理.pdf，SHA-256 340aba76699a12f330637f9c682abbc893e03e0e91788287934c0a7e0a3cace9。
% 原页：第 1 页（A3 六栏）第 2 栏下部，书页码 635。
% 重绘依据：正视图、侧视图均为宽 1、高 2 的矩形，上方接底宽 2、高 1 的三角形；
%   俯视图为边长 2 的正方形，内含边长 1 的虚线正方形（棱柱顶面），正方形对角线为实线。
\begin{tikzpicture}[semithick, every node/.style={inner sep=1.2pt, font=\small}]
  % 正视图
  \draw (0,4) rectangle (1,6);
  \draw (-0.5,6) -- (0.5,7) -- (1.5,6);
  \draw (-0.5,6) -- (1.5,6);
  \node[below] at (0.5,3.25) {正视图};
  \draw[<->] (-0.9,4) -- (-0.9,6) node[midway,left] {$2$};
  \draw[<->] (-0.9,6) -- (-0.9,7) node[midway,left] {$1$};
  \draw[<->] (0,3.78) -- (1,3.78) node[midway,below] {$1$};
  % 侧视图
  \draw (3,4) rectangle (4,6);
  \draw (2.5,6) -- (3.5,7) -- (4.5,6);
  \draw (2.5,6) -- (4.5,6);
  \node[below] at (3.5,3.25) {侧视图};
  \draw[<->] (2.6,4) -- (2.6,6) node[midway,left] {$2$};
  \draw[<->] (2.6,6) -- (2.6,7) node[midway,left] {$1$};
  \draw[<->] (3,3.78) -- (4,3.78) node[midway,below] {$1$};
  % 俯视图
  \draw (0.5,0) rectangle (2.5,2);
  \draw[dashed] (1,0.5) rectangle (2,1.5);
  \draw (0.5,0) -- (2.5,2);
  \draw (0.5,2) -- (2.5,0);
  \draw[<->] (0.5,-0.38) -- (2.5,-0.38) node[midway,below] {$2$};
  \draw[<->] (0.1,0) -- (0.1,2) node[midway,left] {$2$};
  \node[below] at (1.5,-0.9) {俯视图};
\end{tikzpicture}
